\cotefile{29}
\scanpage{29}
% archive sheet number: 28

est``l'image inverse'' de $\underline{R}_{(S)}$, toujours sous les conditions de l'exemple 3. Alors il est \uline{trivial} que $\operatorname{Rev}^{\underline{R}_{(S)}}(S)\qquad\operatorname{Rev}^{\underline{R}_{(T)}}(T)$ est une équivalence, utilisant le résultat analogue pour les revêtements étales de schémas $X$ et $X\unsure{x|\times}_{S}T$ (où $X\ \deleted{\unsure{\mathrm{est}}}=\underline{s}(M)$ est un schéma \deleted{\unsure{cxxvx|connx}} sur $S$ ..). En particulier, quand on dispose d'un point géométrique non ramifié de $T$, on obtient un isomorphisme sur les groupes fondamentaux. Le même raisonnement, après localisation étale à peu près évidente, s'applique au\deleted{\unsure{x|s}} cas kummérien \deleted{\unsure{a)}}.

J'ai fait allusion p\deleted{\unsure{h}}lusieurs fois à la donnée de définir une donnée de ramification image inverse d'une donnée de ramification sur $S$ par un morphisme $T\quad S$; la définition explicite serait nécessairement un peu lourde; je me borne à signaler que dans le cas où $\underline{R}$ est défini par $(Z,G)$ comme dans l'exemple 1, son image inverse est définie par $(Z_T=Zx_ST,G)$. L'ennui avec cette définition générale, c'est qu'elle ne s'applique pas directement au cas où $\underline{R}$ est une donnée kummérienne associée à un systême $(D_i)$, car la donnée image inverse au sens précis précédent n'\uline{est pas} la donnée correspondante au\deleted{\unsure{xxRx|xxex}} systême $(D_{i\,T})$ même à équivalence de catégories fibrées près; elle lui est seulement ``essentiellement équivalente'' au sens du n°1. Seule une rédaction en forme permettra de se rendre compte s'il y a lieu d'introduire cette notion d'image inverse, ou s'il ne vaut pas mieux donner des conditions axiomatiques sur la donnée $(\unsure{*|\pi})$, sous les conditions de l'exemple 3, qui assurent moralement que $\underline{R}_{(T)}$ est essentiellement équivalente à l'image inverse de $\underline{R}_{(S)}$. Bien sûr, ces conditions, pour \deleted{\unsure{avoir|avxix}} être utiles, doivent s'appliquer directement aux cas kummériens. Ce sera sans doute le parti le plus raisonnable. Il me semble également qu' serait raisonnable, assez rapidement après les définitions général-
\scanpage{31}
% archive sheet number: 30

passage au complété (par exemple si $\underline{O}_{S,s}$ à fibres formelles géom. normales, par exemple s'il est excellent). Les mêmes observations s'appliqueront sans hypothèse de normalité ou de réduction quelconque; à condition de se borner aux revêtements \deleted{kummériens}\inserted{mod.ram.} \uline{galoisiens} ; cf n° suivant.

\unsure{3|4}. \uline{Cas des revêtements ga\deleted{\unsure{h}}loisiens}.

Un revêtement \uline{galoisien} de type $\underline{R}$ de $S$, à groupe de Galois $G$, est un objet $X$ de $\operatorname{Rev}^{\underline{R}}(S)$, sur lequel le groupe fini $G$ opère à \deleted{\unsure{x}}droite, de tel façon que dans la catégorie $\operatorname{Rev}^{\underline{R}}(S)$ on obtienne un\deleted{\unsure{x|e}} revêtement principal de l'objet final, i.e. 1°) le morphisme $X\unsure{x|\times}G\quad XxX$ est un isomorphisme et 2°) le morphisme $X\quad\text{objet final}$ est couvrant. Comme $\operatorname{Rev}^{\underline{R}}(S)$ est en tous cas une catégorie multigaloisienne, la notion de revêtement galoisi\unsure{\.{e}|e|é}n de type $\underline{R}$ s'explicite également en termes des groupesfondamentaux décrivant cette catégorie. On peut enfin l'expliciter en revenant à la définition de $\operatorname{Rev}^{\underline{R}}(S)$, et en exigeant que les revêtements des $\underline{s}(M)$ $(M\quad\operatorname{Ob}\underline{C}_{S'}\ )$ correspondants à $X$ soient principaux de groupe $G$. Par abus de langage, l\deleted{\unsure{'}}e\deleted{objet}\inserted{revêtement} à opérateurs $(\underline{r}(X),G)$ de $S$ s'appellera également revêtement galoisien de type $\underline{R}$ de $S$. Cet abus de langage ne présente pas d'inconvénient trop sérieux dans le cas où le foncteur de la catégorie des revêtements galoisiens de type $\underline{R}$ de $S$, dans la catégorie des revêtements à opérateurs de $S$, défini par le foncteur ``réalisation géométrique'', est pleinement fidèle,lorsqu'on se borne à prendre comme morphismes dans la première les seuls isomor\deleted{\unsure{phx}}phismes. (Alors, en effet, pour un revêtement à opérateurs $(T,G)$ donné de $S$, il existe, à isomorphisme canonique près, au plus un $(X,G)$ \deleted{dont} galoisien de type $\underline{R}$ dont il provient). Nous allons donner des conditi\deleted{\unsure{o|n}}ons pour qu'il en soit bien ainsi, en nous plaçant comme d'habitude dans le cas de l'exemple 3. Voici d'abord la condition néc. et suffisan
\scanpage{32}
% archive sheet number: 31

te dans le cas particulier de l'exemple 1: pour tout localise strict $S'$ d'un point de $S$, prenant une composante connexe $Z'_o$ de $Z'=Z\unsure{x|\times}_{\unsure{\substack{S\\S}|\mathbb{S}}}S'$ et son stabilisateur $H$ dans $G$, et pour deux sous-groupes \inserted{distingués} $H_1,H_2$ de $H$ et tout i\unsure{d|s}\unsure{\deleted{h}o|o}morphisme $u:Z'_o/H_1^{\deleted{\unsure{1|I}}}\quad Z'_2/H_2$, \deleted{\unsure{ilxx|il}} compatible avec un isomorphisme $\handwritten{\phi}:\unsure{\mathbb{H}|H}/H_1\quad H/H_2$, il existe un élément $g$ de $\mathbb{H}$ tel que $\operatorname{int}(g)$ applique $H_1$ sur $H_2$ et induise $\handwritten{\phi}$ par passage au quotient, et tel que $u$ soit induit par l'endomorphisme $g_{Z'_o}$ de $Z'_o$ par passage aux quotients. Dans le cas général de l'exemple 3, il faut exiger ceci pour tous les couples $(\underline{s}(M),G)$ sur les divers $S'\quad\underline{\operatorname{Et}}_S$.

This condition is satisfied in the case of Kummer data, as one sees by combining 1.10. and 1.15. of Murre's notes. This phenomenon explains why I have been so reluctant (wrongly so) to define \deleted{Kummer}\inserted{tamely ramified} coverings except in the Galois case: in the non galois case, it will have to be not just a covering satisfying some extra \uline{condition\unsure{s|g}}, but a covering with some extra \uline{structure} ; whereas in the Galois case, the extra structure is uniquely determined, if it exists. Note however that \deleted{\unsure{forxmorphism|for morphism}} not all morphisms \inserted{(in the usual naive sense)} between (even Galois) tame coverings which are not isomorphisms should be considered, only those which come form morphisms in $\operatorname{Rev}^{\underline{R}}(S)$, - that one has to be careful is shown by the counterexample to prop.1.16. of Murre's notes. Also, the natural notion of products of tame \deleted{Galois} coverings, whether \unsure{galois|Galois} or not, is not the naïve fibered product over $S$ (which by the way accounts for the complication of the proof of Murre's lemma 4.3.). This shows rather clearly that even in the Galois case, the consideration of the category $\operatorname{Rev}^{\underline{R}}(S)$ is indispensable for a good comprehension of what is really going on.
\scanpage{33}
% archive sheet number: 32

\uline{Remords} concernant l'exemple 3 du par.1. - La présentation de cet exemple n'est ni élégante ni naturelle, la donnée de $\underline{s}'$ et de \quad étant un élément de structure artificiel. Il vaut mieux présenter les choses ainsi, en termes de conditions sur $\underline{R}$, savoir les suivantes:

a) Pour tout objet $M$ d'un $\underline{C}_{S'}$, le faisceau \inserted{étale} associé au préfaisceau sur $\underline{\operatorname{Et}}_{S'}$ des automorphismes de $M$ est \deleted{\unsure{un|du}} représentable par un schéma en groupes $\underline{G}(M)$ fini et étale sur $S$ (évidemment déterminé à isomorphisme unique près.

Il est évident que $G(M)$ opère sur le schéma relatif $\underline{s}(M)\unsure{=|\equiv}Z$.

b) Pour tout objet $S'$ de $\underline{\operatorname{Et}}_S$, tout morphisme \deleted{$M'\unsure{xx}M$} dans $\underline{C}_{S'}$ est un épimorphisme, et tout endomorphisme d'un obje\deleted{\unsure{y|j}}t de $\underline{C}_{S'}$ est un automorphisme. Pour tout morphisme $\inserted{u:}M'\quad M$ dans $\underline{C}_{S'}$, et tout automorphisme $g'$ de $M'$, il existe un automorphisme $g$ de $M$ tel que $gu=ug'$.
% The cancelled M'xxM is typed above the line, above "dans".

Cet automorphisme est évidemment unique, de plus la condition précédente implique que à $\inserted{u:}M'\quad M$ est associé un unique homomorphisme \unsure{de} schémas en groupes $G(M')\quad G(M)$ tel que $u$ soit compatible avec les opér\unsure{a}tions de $G(M')$ et $G(M)$ sur $M'$ et $M$.

c) Avec les notations précédentes, \deleted{si} l'homomorphisme $G(M')\quad G(M)$ est surjectif, et si $H$ désigne son noyau, l\deleted{\unsure{'}}e\deleted{homo}morphisme $\underline{s}(M')\quad\underline{s}(M)$ induit un \deleted{morphisme} isomorphisme $\underline{s}(M')/H\quad\underline{s}(M)$, \deleted{\unsure{i.e.|il} fait de $\underline{s}(\unsure{M'|M})$ un \unsure{fibéé|fibré} principal sur $\underline{s}(M)$ de groupe} (NB on suppose que les $\underline{s}(M)$ sont \deleted{\unsure{relativement} des \unsure{revêtem|revêteme}} relativement affines, de sorte que le passage au quotient a un sens).

d) Deux homomorphismes de $M'$\deleted{\unsure{et|à}} dans $M$ sont,localement pour la top\unsure{o|\deleted{t}o|\deleted{I}o}logie étale, déduits l'un de l'autre par composition avec un \unsure{automorphisme} de $M$.
\scanpage{35}
% archive sheet number: 34

Dear Murre,\\

I am very sorry I did not succeed to convey the intuitive idea behin\inserted{d} the general nonsense of my notes. It seems to me that the basic example in order to understand the idea is example 1, where you can take $Z$ to be a standard Kummer covering for definiteness, $Z=Z_a^n$, and $S$ normal. Intuitively, when you look at coverings \inserted{(normal, say) $S'$} of $S$ whose ramification type is not worse than the one of $Z$ over $S$, you mean that the \inserted{normalized} inverse image $Z'$ of $S'$ over $Z$ is étale. \deleted{But by transport de} From the birational point of view, assuming $S'$ connected and therefore corresponding to a field extension $K'$ of the field of functions $K$ of $S$, this means simply that $K'$ is isomorphic to a sub\inserted{-$K$-}extension of an extension \inserted{$L'$} of the function field $L$ of $Z$, unramified with respect to the model $Z$; when $S$ (and whence $Z$) is strictly local, \inserted{hence $L'=L$,} this means simply that $K'$ is isomorhic to a subextension of $L$, a very reasonable way indeed of interpreting the intuitive fact that $S'/S$ has no worse ramification than $Z/S$. (NB We assume of course now,for simplicity, that $Z$ is connected, and mor\unsure{\.{e}|é|e}over that $Z$ is the normalisation of $S$ in some finite Galois extension $L$ of $K$). Note that anyhow, we want that the notion in question should be of local nature for the étale topology, thus we can reduce always to the s rictly local case. Now by transport de structure the group $G$ operating on $Z$ must operate on $S'$\inserted{(in a way compatible with\deleted{\unsure{h|x}} operationson $Z$)}, and it is immediate that we recover $S'$ from the $G$-covering $Z'$ of $Z$ as being simply $S'=Z'/G$. Thus \deleted{\unsure{nexre|hence|here}} the category of coverings of $S$ we are interested in is just equivalent to the catégory of \inserted{étale} coverings of $(Z,G)$. This latter catégory on the other hand is of a very standard type; the presence of $G$ does not prevent it to share all the properties of catégories of étale coverings
